% ===== PRÉAMBULE CANONIQUE — à coller VERBATIM en tête de CHAQUE .tex =====
\documentclass[11pt,a4paper]{article}
\usepackage[utf8]{inputenc}\usepackage[T1]{fontenc}\usepackage{lmodern}
\usepackage[french]{babel}
\usepackage{amsmath,amssymb,mathtools}\usepackage{icomma}
\usepackage[version=4]{mhchem}          % \ce{...} : équations & formules
\usepackage{siunitx}\sisetup{locale=FR} % \qty{}{}, \si{}, \ang{}
\usepackage{chemfig}                    % \chemfig{...} : structures développées
\usepackage{xcolor}\usepackage{colortbl}
\usepackage{tikz}\usepackage{pgfplots}\pgfplotsset{compat=1.18}
\usepackage[most]{tcolorbox}\usepackage{enumitem}\usepackage{array,booktabs}
\usepackage[a4paper,margin=2.2cm]{geometry}
\usetikzlibrary{arrows.meta,calc,angles,quotes,positioning,patterns,backgrounds,shapes.geometric}
\definecolor{primary}{RGB}{222,49,99}\definecolor{primarydark}{RGB}{150,22,55}
\definecolor{defcol}{RGB}{0,114,178}\definecolor{methcol}{RGB}{52,92,148}
\definecolor{excol}{RGB}{0,135,90}\definecolor{attncol}{RGB}{200,40,55}
\tcbset{boxbase/.style={enhanced,breakable,boxrule=0.9pt,arc=2.5pt,
  left=3mm,right=3mm,top=2.3mm,bottom=2.3mm,fonttitle=\bfseries,coltitle=white}}
% --- boîtes TUTORIEL (noms EXACTS, ne pas renommer) ---
\newtcolorbox{cle}[1][]{boxbase,colback=defcol!6,colframe=defcol,
  title={\textsc{À retenir}\ifx\relax#1\relax\relax\else\textmd{ : \itshape #1}\fi}}
\newtcolorbox{methode}[1][]{boxbase,colback=methcol!7,colframe=methcol,sharp corners,
  title={\textsc{Méthode}\ifx\relax#1\relax\relax\else\textmd{ --- \itshape #1}\fi}}
\newtcolorbox{exemplebox}[1][]{boxbase,colback=excol!5,colframe=excol,
  title={\textsc{Exemple}\ifx\relax#1\relax\relax\else\textmd{ : \itshape #1}\fi}}
\newtcolorbox{piege}[1][]{boxbase,colback=attncol!6,colframe=attncol,
  title={\textsc{Piège}\ifx\relax#1\relax\relax\else\textmd{ : \itshape #1}\fi}}
\newtcolorbox{remarque}[1][]{boxbase,colback=black!4,colframe=black!45,
  title={\textsc{Remarque}\ifx\relax#1\relax\relax\else\textmd{ : \itshape #1}\fi}}
% --- boîtes EXERCICE / CORRIGÉ (noms EXACTS, auto-comptées) ---
\newtcolorbox[auto counter]{exo}[1][]{boxbase,colback=primary!4,colframe=primary,
  title={\textsc{Exercice}~\thetcbcounter\ifx\relax#1\relax\relax\else\textmd{ --- \itshape #1}\fi}}
\newtcolorbox[auto counter]{corrige}[1][]{boxbase,colback=excol!4,colframe=excol,
  title={\textsc{Corrigé}~\thetcbcounter\ifx\relax#1\relax\relax\else\textmd{ --- \itshape #1}\fi}}
\newcommand{\R}{\mathbb{R}}\newcommand{\N}{\mathbb{N}}\newcommand{\Z}{\mathbb{Z}}\newcommand{\ds}{\displaystyle}
% (un \usepackage{fancyhdr}+en-tête est possible ; il est ignoré côté web)
% =========================================================================
\begin{document}

\begin{center}{\large\bfseries\color{primarydark} Thème 1 — Corrigé Moyen}\end{center}

\begin{corrige}[chaîne masse $\to$ moles $\to$ molécules]
On enchaîne $m \xrightarrow{\div M} n \xrightarrow{\times N_{\mathrm{A}}} N$.
\begin{enumerate}[label=\alph*)]
  \item $n = \dfrac{9}{18} = \SI{0.5}{\mole}$, puis $N = 0{,}5\times 6{,}02\times10^{23}
        = 3{,}01\times10^{23}$ molécules.
  \item $n = \dfrac{8{,}8}{44} = \SI{0.2}{\mole}$, puis $N = 0{,}2\times 6{,}02\times10^{23}
        = 1{,}20\times10^{23}$ molécules.
\end{enumerate}
\end{corrige}

\begin{corrige}[masse molaire du glucose, puis quantité de matière]
\begin{enumerate}[label=\alph*)]
  \item $M(\ce{C6H12O6}) = 6\times 12 + 12\times 1 + 6\times 16 = 72+12+96 = \SI{180}{\gram\per\mole}$.
  \item $n = \dfrac{90}{180} = \SI{0.500}{\mole}$.
  \item $n = \dfrac{18}{180} = \SI{0.100}{\mole}$.
\end{enumerate}
\end{corrige}

\begin{corrige}[masse d'un nombre d'entités donné]
On enchaîne $N \xrightarrow{\div N_{\mathrm{A}}} n \xrightarrow{\times M} m$.
\begin{enumerate}[label=\alph*)]
  \item $n = \dfrac{3{,}01\times10^{23}}{6{,}02\times10^{23}} = \SI{0.5}{\mole}$, puis
        $m = 0{,}5\times 32 = \SI{16.0}{\gram}$.
  \item $n = \dfrac{1{,}204\times10^{24}}{6{,}02\times10^{23}} = \SI{2}{\mole}$, puis
        $m = 2\times 56 = \SI{112}{\gram}$.
\end{enumerate}
\end{corrige}

\begin{corrige}[pourcentage massique d'un élément]
Le pourcentage massique d'un élément vaut $\dfrac{\text{masse de l'élément dans 1 mol}}{M}\times 100$.
\begin{enumerate}[label=\alph*)]
  \item $\%\,\ce{O} = \dfrac{16}{18}\times 100 \approx 88{,}9\,\%$.
  \item L'azote apparaît \emph{deux} fois : masse $2\times 14 = 28$. $\%\,\ce{N} = \dfrac{28}{80}\times 100 = 35{,}0\,\%$.
  \item $\%\,\ce{Ca} = \dfrac{40}{100}\times 100 = 40{,}0\,\%$.
\end{enumerate}
\end{corrige}

\begin{corrige}[compter les atomes]
\begin{enumerate}[label=\alph*)]
  \item \ce{H2SO4} porte $4$ atomes \ce{O} : $n(\ce{O}) = 0{,}2\times 4 = \SI{0.8}{\mole}$, soit
        $N = 0{,}8\times 6{,}02\times10^{23} = 4{,}82\times10^{23}$ atomes d'oxygène.
  \item \ce{O2} : $n = \dfrac{8}{32} = \SI{0.25}{\mole}$ de \emph{molécules}, soit
        $0{,}25\times 2 = \SI{0.5}{\mole}$ d'atomes. \ce{S} : $n = \dfrac{8}{32} = \SI{0.25}{\mole}$ d'atomes.
        À masse égale, \ce{O2} contient donc \textbf{deux fois plus} d'atomes que le soufre.
\end{enumerate}
\end{corrige}

\end{document}
